\section{What the Environmental Comparisons Establish}
\label{sec:environmental-evidence}
Part~II is potentially more informative than another fit to the same
kinship moments: it asks how outcomes change with family circumstances.
But a finding supports genetic transmission over an environmental account
only if the competing accounts predict different findings under the
conditions studied. A null estimate for one exposure is not a test of all
social mechanisms. The intervention, population, outcome, and uncertainty
must remain explicit.

\subsection{Regression Towards the Mean Does Not Identify Its Cause}
\label{sec:regression-mean}

Chapter~14 documents approximately linear relationships between parental
outcomes and mean offspring outcomes for occupation, longevity, log house
value and neighbourhood deprivation. The plotted slopes are below one:
children of parents at opposite extremes have less extreme average
outcomes. This is a genuine descriptive pattern of regression towards the
mean, comparable in form to Galton's height example. The nearly flat
longevity relationship indicates particularly strong regression, not its
absence. The high reported $R^2$ values describe the fit to grouped means,
not the share of individual offspring variation explained by parents
\citep[pp.~210--216]{Clark2026}.

The difficulty is the inference from this pattern to its mechanism.
Clark calls it ``confirmation of the likely dominance of direct genetic
transmission in determining most social outcomes'' (p.~216). Yet a purely
cultural model can generate the same linear, symmetric conditional mean.
For centred parental and offspring cultural characteristics, let
\[
 C_o=\gamma\frac{C_f+C_m}{2}+U,\qquad
 0<\gamma<1,\qquad \mathbb E[U\mid C_f,C_m]=0.
\]
Writing $\overline C=(C_f+C_m)/2$, the model implies
$\mathbb E[C_o\mid\overline C]=\gamma\overline C$: deviations of either
sign contract by the same proportion, without genetic transmission.
Cultural-transmission models with this property are also developed by
\citet{OttoEtAl1994}. Conversely, genetic additivity alone does not establish an exactly linear
conditional mean; additional distributional or conditional-mean assumptions
are required.

The figures may challenge cultural models predicting strong poverty or
privilege traps. But broad bins and high $R^2$ values do not replace
formal tests of those predictions. Distinguishing genetic from cultural
transmission requires restrictions or evidence on which the mechanisms
differ.


\subsection{Family Circumstances, Compensation, and the Relevant Contrast}
Parental loss illustrates the distinction. Clark argues that predominantly
genetic determination predicts little harm from early parental death
\citep[p.~189]{Clark2026}. Yet he also recognizes the alternative:
``Either nurture was not of great impact, or social mechanisms existed to
replace lost parents'' \citep[p.~288]{Clark2026}. If relatives or public
institutions replace an important input, a small effect of bereavement
need not imply that the input was unimportant. The comparison identifies
the consequences of loss with the responses it induces, not the effect
of removing care without replacement.

Family size and birth order require similarly specific contrasts \citep[ch.~8]{Clark2026}; see also \citet{BlackDevereuxSalvanes2005} for related economic evidence. A
twinning-based comparison must justify why the additional birth isolates
the relevant resource change, who enters the sample, and which families'
responses it identifies. A small effect along that margin does not show
that every aspect of family resources is irrelevant. Nor does parental
symmetry by itself distinguish genetic transmission from social mechanisms
that operate through both parents. These tests should be judged against
specified alternatives rather than a single undifferentiated cultural model.

Exogenous wealth shocks offer a different source of evidence. The Georgia
land-lottery study of \citet{BleakleyFerrie2016}, for example, studies random
wealth and outcomes across generations. Such designs belong in the
assessment because they change resources without inferring that change
from kinship resemblance. Their implications remain specific to the
resource shock, institutions, and generations observed; they do not by
themselves settle the effects of education or other family inputs.

\subsection{Schooling Effects and Within-Twin Comparisons}
An effect of an additional year of schooling is not a heritability estimate,
and neither is the return on an additional pound of educational spending.
Evidence on compulsory-schooling reforms concerns the students induced to
remain in school by a particular reform and the outcomes measured afterward.
Its relevance to other margins of education requires an argument. Likewise,
failure to change intergenerational rank persistence does not establish a
zero effect on earnings levels, health, or well-being.


A related, longer-standing literature makes the same kind of point about
within-family comparisons more generally. MZ-twin-difference designs, widely
used to estimate the return to schooling on the logic that identical twins
share both genotype and family background \citep{Griliches1979}, are only as
good as the assumption that within-pair schooling differences are unrelated
to individual ability \citep{Griliches1977,BoundSolon1999,Neumark1999}.
Directly testing this assumption, \citet{SandewallEtAl2014} find that
within-pair differences in adolescent IQ predict within-pair differences in
educational attainment among monozygotic twins, and that controlling for IQ
reduces the within-pair estimated return to schooling by roughly 15\%. These findings concern the identification of a schooling effect,
not the genetic variance of schooling or earnings. Clark is right that
within-MZ comparisons need not be causal, but the limitation is well
established. Broader work on sibling comparisons likewise examines
nonshared confounding, measurement error, and the assumptions required
for causal interpretation \citep{SjolanderEtAl2012,SjolanderEtAl2022}.
The appropriate response is to assess a design's identifying assumptions,
not infer that schooling has no causal effect.


These distinctions organize the assessment of Clark's intervention-specific
evidence. The model in Part~I does not supply a missing causal estimate in
Part~II, and a particular intervention's result does not validate all of
the exclusions in Part~I. Section~\ref{sec:interpretation} returns to the
further step from causal effects to claims about optimal spending.
